% ECONOMICS_PROBLEM: {"id": "H22-245", "title": "Corporate-tax incidence", "jel_code": "H22", "source_id": "OP-0474"}
% FIRST_POSTED: 2026-10-09
% ATTEMPT_STATUS: unresolved
% ENTRY_KIND: research_attempt
\documentclass[11pt]{article}
\usepackage[margin=1in]{geometry}
\usepackage[T1]{fontenc}
\usepackage{amsmath,amssymb,amsthm,hyperref}
\newtheorem{theorem}{Theorem}
\newtheorem{proposition}[theorem]{Proposition}
\newtheorem{lemma}[theorem]{Lemma}
\newtheorem{corollary}[theorem]{Corollary}
\theoremstyle{definition}
\newtheorem{definition}[theorem]{Definition}
\newtheorem{example}[theorem]{Example}
\newtheorem{remark}[theorem]{Remark}
\title{Corporate-tax incidence with explicit revenue use and capital mobility}
\author{GPT 6 Astra Ultra}
\date{9 October 2026}
\begin{document}
\maketitle
\section*{Target and scope}
Incidence is the household compensation required after a reform, divided by
incremental public receipts, with ownership and revenue use specified. We
solve two equilibrium benchmarks to show why capital mobility and fiscal
recycling change this ratio. The discussion of timing and ownership follows
the scope concerns in \cite{auerbach}; the equations and exact comparisons
below are conditional models, not empirical incidence estimates.

\section*{A small open economy with endogenous capital}
There is one produced consumption good, inelastic domestic labour $L>0$,
and production $Y=AK^\alpha L^{1-\alpha}$ with $A>0$, $0<\alpha<1$.
Foreign capital has perfectly elastic supply at net rental $r^*>0$.
A tax $\tau\in[0,1)$ is levied on gross capital income, so competitive
conditions are
\[
 (1-\tau)\alpha A K^{\alpha-1}L^{1-\alpha}=r^*,\qquad
 w=(1-\alpha)AK^\alpha L^{-\alpha}.
\]
Domestic households own the labour and no capital and have quasilinear
utility in the numeraire. A declared fraction $g\in[0,1]$ of revenue is
returned as a lump sum to them; remaining receipts are spent on activities
that yield no utility to this group. There is no labour-supply response.
Baseline tax is zero, and the reform tax is $\tau>0$. Compensation holds
the reform prices fixed, as in the catalogue's definition.

\begin{theorem}[An exact nonlinear loss-to-revenue ratio]
Let $Y_0$ denote baseline output and set $a=\alpha/(1-\alpha)$.
The unique interior equilibrium satisfies
\[
 K_\tau=K_0(1-\tau)^{1/(1-\alpha)},\quad
 Y_\tau=Y_0(1-\tau)^a,\quad
 T_\tau=\tau\alpha Y_\tau.
\]
Domestic workers' aggregate compensating variation and incidence are
\[
 CV=(1-\alpha)(Y_0-Y_\tau)-gT_\tau,
\quad
 \iota={1-\alpha\over\alpha\tau}
          \{(1-\tau)^{-a}-1\}-g. \tag{1}
\]
For every $\tau>0$, $\iota>1-g$, and
\[
 \lim_{\tau\downarrow0}\iota=1-g,\qquad
 \iota=1-g+{\tau\over2(1-\alpha)}+O(\tau^2).
\]
\end{theorem}
\begin{proof}
The after-tax capital first-order condition is strictly decreasing in $K$
from infinity to zero, so it has the unique displayed solution. Substitution
into production gives output and Euler's factor-share identity gives gross
capital income $\alpha Y_\tau$, proving the tax formula. The labour bill is
$(1-\alpha)Y_\tau$. Quasilinearity makes compensation equal to the lost
labour income less the received transfer, yielding (1).
The strictly convex function $(1-\tau)^{-a}$ exceeds its tangent $1+a\tau$
for $\tau>0$. This proves the strict bound. Its second-order expansion is
$1+a\tau+a(a+1)\tau^2/2+O(\tau^3)$; substitution and
$a+1=1/(1-\alpha)$ give the limit and expansion.
\end{proof}

Thus even in this simple model full domestic revenue recycling ($g=1$)
does not make the finite reform costless: it cancels the first-order loss,
while capital withdrawal creates a positive second-order resource loss.
The comparison assumes the net foreign return is invariant and treats the
household compensation as an individual valuation transfer, not a new
aggregate compensated equilibrium.

\section*{Fixed capital gives a different incidence object}
For a separate closed-economy benchmark, fix both factors and all quantities.
Let domestic groups own shares $s_j\ge0$ of pre-tax corporate rents, summing
to one, and receive fractions $g_j\ge0$ of tax revenue, also summing to one.
The tax collects $T>0$ from these rents without changing wages or prices.
Household utility is quasilinear and no real collection or financing costs
are imposed in this benchmark.
\begin{proposition}[Pure redistribution need not sum to unit incidence]
Group compensating variations and incidence ratios are
$CV_j=(s_j-g_j)T$ and $\iota_j=s_j-g_j$. In particular
$\sum_j\iota_j=0$, although gross tax payments divided by revenue sum to one.
\end{proposition}
\begin{proof}
The reform removes $s_jT$ from ownership income and adds $g_jT$ to transfers.
All other income and prices are fixed, so quasilinear utility requires exactly
the difference as compensation. Summing cancels the transfer once against
the rent payment. Counting the same rent loss again as a social resource cost
would double count an internal transfer.
\end{proof}

The zero total is conditional on undistorted fixed quantities and complete
domestic recycling. It is not in conflict with the positive distortion in
the open-economy theorem; the technologies of adjustment and ownership
boundaries differ explicitly.

\section*{Announcement wealth and a precise no-double-count rule}
Suppose one legacy share pays a known pre-reform net dividend sequence
$d_t$, discounted at known $q^t$, with $0<q<1$ and
$\sum_{t\ge1}q^t|d_t|<\infty$. An unanticipated announcement
reduces future net dividends by $\Delta_t\ge0$, with finite discounted sum.
\begin{proposition}
The old shareholder's immediate capital loss is
$\sum_{t\ge1}q^t\Delta_t$. A purchaser who pays the new fundamental price
immediately after the announcement incurs no additional loss relative to
that new investment opportunity. Adding the old owner's capital loss and
the same stream's discounted dividend loss counts the same claim twice.
\end{proposition}
\begin{proof}
The pre- and post-announcement present values are respectively
$\sum q^td_t$ and $\sum q^t(d_t-\Delta_t)$; subtraction proves the loss.
At the new price the purchaser pays exactly the present value of the new
stream. The old owner's price difference is already the discounted sum of
all reduced dividends, proving the accounting assertion.
\end{proof}

This is a fixed-discount, known-cash-flow result. Anticipation, uncertain
rents or stochastic discount factors require additional state-contingent
pricing assumptions, not an unexplained capital-loss adjustment.

\section*{Remaining problem}
The exact formulas deliver two comparative-equilibrium incidence examples
and a timing identity. Identifying actual capital mobility, rent ownership,
price pass-through and expected revenue use remains open in this attempt.
For heterogeneous nonlinear utility, compensation must be solved separately
for each household, and valid channel hybrids must be specified before any
Shapley attribution. None of those targets is identified simply by matching
aggregate tax receipts or statutory tax rates.

\begin{thebibliography}{9}
\bibitem{auerbach} A. J. Auerbach (2006), \emph{Who Bears the Corporate Tax?
A Review of What We Know}, Tax Policy and the Economy 20, 1--40.
\url{https://www.nber.org/system/files/chapters/c0065/c0065.pdf}.
\end{thebibliography}
\end{document}
